% Slide — Formulação restrita
\begin{frame}{Estágio 2: formulação restrita (Eq.~4.2)}
  \footnotesize
  \begin{columns}[T,onlytextwidth]
    \column{0.5\textwidth}
      \textbf{Objetivo, a minimizar:}
      \[
        \mathrm{custo} = \mathrm{CTI} + \delta \cdot p \cdot
        \max(\mathrm{CTI}, 1)
      \]
      \[
        \delta = \max(0,\ \alpha_{\min} - \mathrm{NS})
      \]
      em que $\alpha_{\min}=0{,}70$ e $p=10$ (peso de penalidade).
    \column{0.5\textwidth}
      \textbf{Por que isso resolve o problema de escala}
      \begin{itemize}
        \item A penalidade é \highlight{proporcional ao próprio CTI}
          (via $\max(\mathrm{CTI},1)$), não uma constante fixa.
        \item Candidatos viáveis ($\mathrm{NS}\geq\alpha_{\min}$) são
          ordenados por $-\mathrm{CTI}$ puro -- sem distorção de escala.
        \item Mesma fórmula usada nas versões escalar
          (\texttt{policies.\_eval\_static}) e vetorizada
          (\texttt{constrained\_cost}) -- \textbf{idênticas por construção}.
      \end{itemize}
  \end{columns}
  \vspace{0.8em}

  \small \textbf{Retrocompatibilidade:} \texttt{fitness\_mode="weighted"}
  preserva o comportamento antigo, para reproduzir os resultados já
  reportados no Capítulo 5 se necessário.

  \note{
    A formulação restrita segue exatamente a Equação 4.2 da proposta:
    minimizar CTI sujeito a NS maior ou igual a 0,70. Em vez de descartar
    candidatos inviáveis (o que perderia informação de gradiente para a
    busca), a restrição entra como penalidade proporcional ao déficit de
    serviço, escalada pelo próprio CTI da trajetória. Isso é o que resolve
    o problema de escala: como a penalidade é relativa, não absoluta, o
    mesmo critério funciona igualmente bem numa loja de CTI=50 e numa de
    CTI=6000. O modo antigo continua disponível por retrocompatibilidade,
    caso seja necessário reproduzir exatamente os números já publicados no
    capítulo de resultados.
  }
\end{frame}
